% =============================================================================
%  Acelerador em FPGA do Supervisor de Seguranca LiDAR de uma Cadeira de Rodas
%  Projeto Final -- PCR / Co-Projeto Hardware-Software -- UnB 2026/1
%
%  Compilacao: pdflatex relatorio_pcr.tex  (2x, por causa das referencias)
%  Requer os .dat em docs/figs/ (gerados por golden/gen_figs.py)
%
%  Numeros de sintese (Vivado 2022.2, xc7z020clg484-1) ja preenchidos:
%    WNS +0,872 ns | Fmax 109,6 MHz | LUT 2445 | FF 1808 | DSP 54 | BRAM 0
%    Potencia 0,246 W (0,140 dinamica + 0,106 estatica)
%  Reproduzir: vivado -mode batch -source sim/report_only.tcl
% =============================================================================
\documentclass[conference]{IEEEtran}
\IEEEoverridecommandlockouts

\usepackage[utf8]{inputenc}
\usepackage[T1]{fontenc}
\usepackage[brazilian]{babel}
\usepackage{graphicx}
\usepackage{amsmath,amssymb}
\usepackage{booktabs}
\usepackage{cite}
\usepackage{url}
\usepackage{xcolor}
\usepackage{tikz}
\usepackage{pgfplots}
\usepackage{listings}
\usepackage{siunitx}
\pgfplotsset{compat=1.17}
\usetikzlibrary{arrows.meta, positioning, fit, calc}

\sisetup{output-decimal-marker={,}, per-mode=symbol, detect-all}

% Marcava numeros pendentes do Vivado. Nao ha mais nenhum -- mantida por
% compatibilidade caso se acrescente algum campo novo.
\newcommand{\vivado}[1]{\textbf{#1}}

\lstset{basicstyle=\ttfamily\scriptsize, breaklines=true, columns=fullflexible,
        frame=single, framesep=3pt, xleftmargin=2pt}

\begin{document}

\title{Acelerador em FPGA para o Supervisor de Seguran\c{c}a\\
LiDAR de uma Cadeira de Rodas Semi-Assistida:\\
Co-Projeto Hardware--Software sobre Zynq-7000}

\author{
\IEEEauthorblockN{Vítor Gonçalves de Andrade Silva}
\IEEEauthorblockA{Engenharia Eletrônica --- FCTE/UnB\\
Email: 261025182@aluno.unb.br}
\and
\IEEEauthorblockN{Thiago Henrique Oliveira Souza}
\IEEEauthorblockA{Engenharia Eletrônica --- FCTE/UnB\\
Email: 180131672@aluno.unb.br}
}

\maketitle

% =============================================================================
\begin{abstract}
Este trabalho acelera em FPGA o n\'ucleo do supervisor de seguran\c{c}a de uma
cadeira de rodas motorizada semi-assistida. O supervisor decide, a cada varredura
de LiDAR, se o avan\c{c}o comandado pelo usu\'ario deve ser reduzido ou
interrompido: \'e o bloco de maior custo do sistema e o \'unico cuja lat\^encia
tem consequ\^encia f\'isica direta, pois determina a dist\^ancia percorrida antes
da frenagem. O algoritmo projeta os \num{720} feixes da varredura no referencial
do ve\'iculo e testa, para \num{19} dire\c{c}\~oes candidatas de desvio, se cada
ponto cai no corredor que a cadeira ocupar\'a --- \num{9918} avalia\c{c}\~oes por
varredura. Prop\~oe-se um particionamento em que o custo $O(N\cdot K)$ migra para
hardware e o custo $O(K)$ (fun\c{c}\~ao custo e escolha do candidato) permanece no
ARM, onde residem os ganhos ajust\'aveis em campo. A arquitetura emprega um CORDIC
pipelinado de \num{16} est\'agios, que n\~ao consome multiplicadores, alimentando
\num{19} pistas paralelas cujas rota\c{c}\~oes t\^em coeficientes constantes de
s\'intese. A comunica\c{c}\~ao usa AXI4-Stream para os feixes (via AXI-DMA) e
AXI4-Lite para calibra\c{c}\~ao e resultados. A equival\^encia num\'erica foi
verificada bit a bit contra um modelo de refer\^encia em ponto fixo, sobre tr\^es
varreduras \emph{reais} capturadas na cadeira; as tr\^es passam. A an\'alise de
precis\~ao revelou um defeito real de formato: o \^angulo do feixe atinge
\SI{-5.85}{\radian} e n\~ao cabe em Q3.13, saturando e projetando o ponto no lugar
errado (erro de \SI{402}{\milli\metre}); com o \^angulo envolvido em
$[-\pi,\pi)$ o erro cai a \SI{0.99}{\milli\metre} m\'edio, cerca de um LSB. Medido
na Raspberry~Pi~5, o software gasta \SI{3.54}{\milli\second} (escalar) ou
\SI{0.488}{\milli\second} (vetorizado); o hardware conclui em \num{751} ciclos
(\SI{7.51}{\micro\second} a \SI{100}{\mega\hertz}), \num{472}$\times$ e
\num{65}$\times$ mais r\'apido, respectivamente.
\end{abstract}

% =============================================================================
\section{Introdu\c{c}\~ao}

Cadeiras de rodas motorizadas s\~ao inacess\'iveis a uma parcela expressiva dos
usu\'arios que mais precisariam delas: pesquisa cl\'inica cl\'assica aponta que
profissionais de reabilita\c{c}\~ao consideram muitos pacientes inaptos ao uso
justamente pela inadequa\c{c}\~ao das interfaces de controle~\cite{fehr2000}.
Espasmos e comandos acidentais no joystick produzem acelera\c{c}\~oes repentinas
e colis\~oes.

O projeto no qual este trabalho se insere converte uma cadeira comercial em um
ve\'iculo semi-assistido sob um princ\'ipio de seguran\c{c}a expl\'icito: o
sistema de alto n\'ivel pode \textbf{frear ou desviar}, mas \textbf{nunca
acelerar}; a iniciativa de movimento \'e sempre do usu\'ario. Um supervisor
executando em Raspberry~Pi~5 com ROS\,2~\cite{ros2docs} l\^e um LiDAR RPLIDAR~C1 e
atenua o comando do joystick antes de repass\'a-lo ao firmware de tempo real em
ESP32-S3.

\subsection{O bloco a acelerar e por qu\^e}

O supervisor \'e, disparadamente, o n\'o mais pesado do sistema: medimos
\SI{17.7}{\percent} da CPU da Pi~5 na implementa\c{c}\~ao escalar e
\SI{8.3}{\percent} ap\'os vetoriza\c{c}\~ao. Mais importante que o custo, por\'em,
\'e a natureza da sua lat\^encia: \emph{ela tem consequ\^encia f\'isica}. O
supervisor s\'o corta o comando de avan\c{c}o; a cadeira ainda percorre uma
dist\^ancia por in\'ercia. Todo atraso entre a chegada da varredura e a decis\~ao
se traduz em cent\'imetros adicionais rumo ao obst\'aculo. Trata-se, portanto, de
um caso em que acelerar n\~ao \'e otimiza\c{c}\~ao, e sim seguran\c{c}a --- o que
motiva o co-projeto.

A escolha \'e ainda mais pertinente porque uma revis\~ao recente do supervisor
revelou uma falha geom\'etrica de seguran\c{c}a. A folga frontal era medida por um
\emph{cone angular} ancorado no sensor, constru\c{c}\~ao v\'alida apenas se o
LiDAR estivesse sobre o eixo longitudinal do ve\'iculo. Estando ele
\SI{0.336}{\metre} fora do eixo, o cone varre uma faixa deslocada e a metade
direita da cadeira fica cega justamente nas curtas dist\^ancias
(Fig.~\ref{fig:cobertura}). A corre\c{c}\~ao substituiu o cone pelo teste do
\emph{corredor f\'isico} que a cadeira ocupar\'a --- e \'e esse algoritmo,
correto por\'em mais caro, que se acelera aqui.

\begin{figure}[!t]
\centering
\begin{tikzpicture}
\begin{axis}[width=\columnwidth, height=4.2cm,
  xlabel={dist\^ancia do obst\'aculo (m)},
  ylabel={cobertura (\%)}, ymin=0, ymax=105, xmin=0.3, xmax=1.5,
  legend pos=south east, legend style={font=\scriptsize},
  grid=major, grid style={gray!20}, tick label style={font=\scriptsize},
  label style={font=\scriptsize}]
  \addplot[thick, red!70!black] table[x=d, y=cone30] {figs/cobertura.dat};
  \addlegendentry{cone $\pm30^\circ$}
  \addplot[thick, orange, dashed] table[x=d, y=cone20] {figs/cobertura.dat};
  \addlegendentry{cone $\pm20^\circ$}
  \addplot[thick, blue!60!black, densely dotted] coordinates {(0.3,100) (1.5,100)};
  \addlegendentry{corredor (proposto)}
\end{axis}
\end{tikzpicture}
\caption{Fra\c{c}\~ao da largura da cadeira efetivamente monitorada. O cone
ancorado no sensor \emph{piora} quanto mais perto est\'a o obst\'aculo, porque
estreita enquanto o desvio lateral do v\'ertice permanece constante: a
\SI{0.60}{\metre} cobre \SI{52}{\percent}, e a parcela n\~ao coberta \'e sempre a
metade direita. O teste de corredor cobre \SI{100}{\percent} por
constru\c{c}\~ao.}
\label{fig:cobertura}
\end{figure}

\subsection{Trabalhos relacionados}

A fun\c{c}\~ao custo que pontua dire\c{c}\~oes candidatas pela folga \'e
aparentada ao \emph{Vector Field Histogram}~\cite{borenstein1991}, com a
distin\c{c}\~ao essencial de raciocinar sobre a \emph{planta} do rob\^o e n\~ao
sobre um cone do sensor. O CORDIC~\cite{volder1959,andraka1998} \'e a escolha
cl\'assica para trigonometria em FPGA sem multiplicadores. O uso de aceleradores
acoplados a processadores embarcados via AXI4 sobre Zynq segue a pr\'atica
consolidada em~\cite{zynqbook}.

% =============================================================================
\section{Desenvolvimento}

\subsection{O algoritmo}

Cada retorno do LiDAR, um par $(r_i, a_i)$, \'e projetado no referencial
\texttt{base\_link} (origem no centro do eixo traseiro):
\begin{align}
X_i &= x_L + r_i\cos(a_i + \theta_L), &
Y_i &= y_L + r_i\sin(a_i + \theta_L),
\label{eq:proj}
\end{align}
\noindent onde $(x_L, y_L, \theta_L)$ \'e a pose calibrada do sensor. Para avaliar
um candidato de desvio $\delta_k$, giram-se os pontos por $-\delta_k$ em torno do
eixo traseiro --- o centro instant\^aneo de rota\c{c}\~ao do acionamento
diferencial:
\begin{align}
x^{(k)}_i &= X_i\cos\delta_k + Y_i\sin\delta_k, \\
y^{(k)}_i &= -X_i\sin\delta_k + Y_i\cos\delta_k .
\label{eq:rot}
\end{align}
O ponto \'e obst\'aculo se $|y^{(k)}_i| \le W/2$ e $x^{(k)}_i > X_f$, com folga
$x^{(k)}_i - X_f$; a folga do candidato \'e o m\'inimo sobre os pontos. Aqui
$W=\SI{0.61}{\metre}$ \'e a largura da cadeira e $X_f=\SI{0.667}{\metre}$ a frente
do chassi. Note que a estrutura da pr\'opria cadeira cai \emph{atr\'as} de $X_f$ e
\'e descartada por constru\c{c}\~ao, sem nenhum limiar arbitr\'ario de
dist\^ancia m\'inima.

Com $N=\num{720}$ feixes e $K=\num{19}$ candidatos ($\delta_k$ de $-45^\circ$ a
$+45^\circ$, passo $5^\circ$), s\~ao \num{9918} avalia\c{c}\~oes por varredura ---
o custo $O(N\cdot K)$ que motiva o acelerador.

\subsection{Particionamento Hardware--Software}

O supervisor decomp\~oe-se em duas partes de custo muito distinto
(Tab.~\ref{tab:part}). Apenas a primeira migra para hardware.

\begin{table}[!t]
\caption{Particionamento HW/SW}
\label{tab:part}
\centering\footnotesize
\begin{tabular}{@{}lll@{}}
\toprule
\textbf{Etapa} & \textbf{Custo} & \textbf{Onde} \\
\midrule
Proje\c{c}\~ao + teste de corredor & $O(N\!\cdot\!K)=\num{9918}$ & \textbf{FPGA} \\
Fun\c{c}\~ao custo + \emph{argmin} & $O(K)=\num{19}$ & ARM \\
\bottomrule
\end{tabular}
\end{table}

A decis\~ao de \emph{n\~ao} levar tudo para hardware \'e deliberada. A fun\c{c}\~ao
custo s\~ao \num{19} compara\c{c}\~oes: o ganho seria irris\'orio. Em compensa\c{c}\~ao,
\'e nela que residem os pesos ajustados em campo ($w_{\text{obs}}$,
$w_{\text{dev}}$, $K_{\text{assist}}$), retocados a cada ensaio. Congel\'a-los em
RTL trocaria flexibilidade por nada --- exatamente o \emph{trade-off} que o
co-projeto existe para arbitrar.

\subsection{Arquitetura de hardware}

A Fig.~\ref{fig:arq} mostra o \emph{datapath}. O CORDIC \'e \textbf{compartilhado}
porque a proje\c{c}\~ao~\eqref{eq:proj} \'e comum a todos os candidatos: o custo
$O(N)$ \'e pago uma \'unica vez. O que se replica s\~ao as \num{19} pistas, que
carregam o custo $O(N\cdot K)$ --- e \'e precisamente a\'i que o paralelismo
espacial da FPGA rende: as pistas avaliam \emph{no mesmo ciclo} aquilo que o
processador precisa iterar.

\begin{figure}[!t]
\centering
\begin{tikzpicture}[font=\scriptsize, node distance=3.5mm,
  blk/.style={draw, rounded corners=2pt, minimum height=7mm, inner sep=3pt,
              fill=gray!8, align=center},
  arr/.style={-Latex}]
  \node[blk] (in) {AXI4-Stream\\$(r_i,\theta_i)$};
  \node[blk, right=of in, fill=blue!8] (cor) {CORDIC\\16 est\'agios\\\textbf{0 DSP}};
  \node[blk, right=of cor] (prj) {proje\c{c}\~ao\\$X_i, Y_i$};
  \node[blk, right=6mm of prj, fill=orange!10, minimum height=13mm]
       (lanes) {19 pistas\\paralelas\\\textbf{76 DSP}};
  \node[blk, below=5mm of lanes] (out) {AXI4-Lite\\19 folgas};
  \draw[arr] (in) -- (cor);
  \draw[arr] (cor) -- node[above, font=\tiny] {$\cos,\sin$} (prj);
  \draw[arr] (prj) -- (lanes);
  \draw[arr] (lanes) -- (out);
\end{tikzpicture}
\caption{\emph{Datapath}. Vaz\~ao de \textbf{1 feixe/ciclo}, sem bolhas. O CORDIC
\'e compartilhado (custo $O(N)$); as pistas s\~ao replicadas (custo $O(N\cdot K)$).}
\label{fig:arq}
\end{figure}

\paragraph{CORDIC}
Optou-se por CORDIC em modo rota\c{c}\~ao~\cite{volder1959}, pipelinado em
\num{16} est\'agios, em vez de uma LUT de seno. Uma LUT exigiria uma ROM indexada
pelo \'indice do feixe, o que s\'o funciona se \texttt{angle\_min}, o incremento e
o \emph{yaw} forem fixos em s\'intese --- mas eles s\~ao \textbf{par\^ametros de
calibra\c{c}\~ao}, reescritos em tempo de execu\c{c}\~ao por AXI4-Lite. O CORDIC
aceita qualquer \^angulo e, por usar apenas somadores e deslocamentos cabeados,
n\~ao consome \emph{nenhum} DSP48. Como o la\c{c}o s\'o converge para
$|z|\le\SI{1.7434}{\radian}$, um est\'agio de redu\c{c}\~ao de quadrante traz
$\theta$ para $[-\pi/2,\pi/2]$ e inverte o sinal na sa\'ida.

\paragraph{Pistas}
Em cada pista, $\cos\delta_k$ e $\sin\delta_k$ s\~ao \textbf{constantes de
s\'intese} (\emph{generics}): a rota\c{c}\~ao~\eqref{eq:rot} vira quatro
multiplica\c{c}\~oes de coeficiente fixo, mapeadas diretamente em DSP48. Da\'i
$19\times4 = \num{76}$ DSP, mais \num{2} da proje\c{c}\~ao: \num{78} de \num{220}
dispon\'iveis no \texttt{xc7z020} (\SI{35}{\percent}).

\subsection{Formato de ponto fixo}

Os formatos adotados est\~ao na Tab.~\ref{tab:fmt}. A escolha da \emph{faixa} de
Q3.13 mostrar-se-\'a decisiva na Se\c{c}\~ao~\ref{sec:prec}.

\begin{table}[!t]
\caption{Formatos adotados}
\label{tab:fmt}
\centering\footnotesize
\begin{tabular}{@{}llll@{}}
\toprule
\textbf{Grandeza} & \textbf{Formato} & \textbf{Faixa} & \textbf{LSB} \\
\midrule
dist\^ancias, coords. & Q6.10 & $\pm\num{32}$\,m & \SI{0.98}{\milli\metre} \\
$\cos$, $\sin$ & Q1.14 & $\pm 2$ & \num{6.1e-5} \\
\^angulos & Q3.13 & $\pm\SI{4.0}{\radian}$ & \SI{1.22e-4}{\radian} \\
\bottomrule
\end{tabular}
\end{table}

\subsection{Interfaces AXI4}

Os feixes chegam por \textbf{AXI4-Stream} (\texttt{TDATA} $=$ [$\theta$(31:16) $|$
$r$(15:0)], \texttt{TUSER} $=$ feixe v\'alido, \texttt{TLAST} $=$ fim da
varredura), alimentado por um AXI-DMA que l\^e a varredura da DDR. O controle e o
resultado usam \textbf{AXI4-Lite}: \texttt{CTRL} (\emph{start}), \texttt{STATUS}
(\emph{done}/\emph{busy}) e \num{19} registradores de folga, com sinal estendido a
\num{32} bits. Uma interrup\c{c}\~ao sinaliza a conclus\~ao, dispensando espera
ocupada no ARM. A escolha \'e natural: \emph{stream} para o fluxo volumoso e
sequencial (a varredura), mapeado em mem\'oria para o punhado de escalares
(calibra\c{c}\~ao e resultado).

% =============================================================================
\section{Resultados}

\subsection{Protocolo de verifica\c{c}\~ao}

A verifica\c{c}\~ao \'e feita por \emph{equival\^encia bit a bit} contra um modelo
de refer\^encia, e n\~ao por inspe\c{c}\~ao de formas de onda. O fluxo \'e:

\begin{enumerate}
\item Um \textbf{modelo em ponto fixo} (Python) emula o RTL exatamente: mesmo
CORDIC, mesmos deslocamentos, mesmos truncamentos.
\item Um \textbf{modelo em ponto flutuante} serve de refer\^encia f\'isica, para
medir o erro de quantiza\c{c}\~ao.
\item O modelo gera o \textbf{est\'imulo} e o \textbf{esperado}; o testbench VHDL
l\^e os arquivos, injeta os \num{720} feixes e compara as \num{19} folgas.
\item As constantes do RTL (\texttt{cordic\_pkg.vhd}) s\~ao \textbf{geradas} pelo
mesmo modelo, de modo que RTL e refer\^encia casam \emph{por constru\c{c}\~ao}, e
n\~ao por transcri\c{c}\~ao manual --- um erro de digita\c{c}\~ao numa tabela de
constantes \'e, aqui, indetect\'avel por inspe\c{c}\~ao.
\end{enumerate}

O est\'imulo n\~ao \'e sint\'etico: s\~ao tr\^es varreduras \textbf{reais} da
cadeira, capturadas com um balde posicionado \SI{21}{\centi\metre} \`a direita do
eixo --- exatamente o ponto cego do supervisor antigo (Fig.~\ref{fig:corredor}).
As tr\^es passam com \textbf{zero diverg\^encia}.

\begin{figure}[!t]
\centering
\begin{tikzpicture}
\begin{axis}[width=\columnwidth, height=5.2cm,
  xlabel={$X$ (m) --- \`a frente do eixo traseiro},
  ylabel={$Y$ (m)}, xmin=0, xmax=2.6, ymin=-1.1, ymax=1.5,
  legend pos=north east, legend style={font=\tiny},
  grid=major, grid style={gray!20}, tick label style={font=\scriptsize},
  label style={font=\scriptsize}, axis equal image=false]
  \addplot[only marks, mark=*, mark size=0.5pt, gray!60]
    table[x=x, y=y] {figs/scan_out.dat};
  \addlegendentry{fora do corredor}
  \addplot[only marks, mark=*, mark size=0.9pt, red!80!black]
    table[x=x, y=y] {figs/scan_in.dat};
  \addlegendentry{obst\'aculo (no corredor)}
  \draw[blue!60!black, thick] (axis cs:0.687,-0.305) rectangle (axis cs:2.6,0.305);
  \node[blue!60!black, font=\tiny] at (axis cs:2.1,0.42) {corredor da cadeira};
  \draw[fill=black] (axis cs:0.667,0.336) circle (1.6pt);
  \node[font=\tiny, anchor=south] at (axis cs:0.667,0.38) {LiDAR};
\end{axis}
\end{tikzpicture}
\caption{Varredura real projetada no \texttt{base\_link} (est\'imulo do
testbench). O LiDAR est\'a \SI{0.336}{\metre} \`a esquerda do eixo. Os
\num{117} pontos em vermelho caem no corredor que a cadeira ocupar\'a --- entre
eles o balde, que o cone antigo n\~ao enxergava.}
\label{fig:corredor}
\end{figure}

\subsection{Precis\~ao: um defeito de formato}
\label{sec:prec}

A an\'alise de precis\~ao produziu o achado mais instrutivo do trabalho. O
primeiro modelo em ponto fixo divergia da refer\^encia em at\'e
\SI{402}{\milli\metre} --- inaceit\'avel num sistema de seguran\c{c}a, e grande
demais para ser arredondamento.

A causa n\~ao era ru\'ido de quantiza\c{c}\~ao, e sim \textbf{estouro de faixa}. O
\^angulo do feixe no referencial do sensor,
$\theta_i = a_{\min} + i\cdot\Delta a + \theta_L$, atinge \SI{-5.85}{\radian}; o
formato Q3.13 comporta apenas $\pm\SI{4.0}{\radian}$. O valor \textbf{saturava}, e
o feixe era projetado numa dire\c{c}\~ao completamente errada. Envolvendo $\theta$
em $[-\pi,\pi)$ antes da quantiza\c{c}\~ao --- opera\c{c}\~ao barata, feita no
host --- o erro desaba (Tab.~\ref{tab:prec}).

\begin{table}[!t]
\caption{Erro do ponto fixo contra a refer\^encia em ponto flutuante}
\label{tab:prec}
\centering\footnotesize
\begin{tabular}{@{}lrr@{}}
\toprule
 & \textbf{erro m\'edio} & \textbf{erro m\'aximo} \\
\midrule
$\theta$ saturado (defeito) & \SI{17.89}{\milli\metre} & \SI{402.15}{\milli\metre} \\
$\theta$ envolvido (corrigido) & \textbf{\SI{0.99}{\milli\metre}} & \textbf{\SI{5.27}{\milli\metre}} \\
\bottomrule
\end{tabular}
\end{table}

O erro final \'e da ordem de \textbf{um LSB} de Q6.10 (\SI{0.98}{\milli\metre}) e
representa \SI{0.9}{\percent} da zona de parada de \SI{600}{\milli\metre} --- ou
seja, a quantiza\c{c}\~ao \emph{n\~ao} \'e o fator limitante do sistema. A
li\c{c}\~ao \'e que, em ponto fixo, a escolha da \textbf{faixa} costuma ser mais
cr\'itica que a da \textbf{resolu\c{c}\~ao}: perdemos \num{402} milímetros por
tr\^es bits de parte inteira, n\~ao por falta de bits fracion\'arios.

A Fig.~\ref{fig:swhw} confronta, candidato a candidato, a folga calculada em
software (ponto flutuante, a refer\^encia f\'isica) com a produzida pelo
\textbf{hardware} --- e o ``hardware'' aqui \'e a sa\'ida efetiva da
simula\c{c}\~ao VHDL, lida de \texttt{rtl\_out}, e n\~ao o modelo. As curvas se
sobrep\~oem: \'e esse o resultado. A Fig.~\ref{fig:erro} mostra o res\'iduo, com
sinal, para as tr\^es varreduras: \num{57} compara\c{c}\~oes, erro m\'edio de
\SI{0.99}{\milli\metre} e m\'aximo de \SI{5.27}{\milli\metre}, todas dentro de
poucos LSB.

\begin{figure}[!t]
\centering
\begin{tikzpicture}
\begin{axis}[width=\columnwidth, height=4.4cm,
  xlabel={candidato de desvio $\delta_k$ (graus)},
  ylabel={folga (m)}, xmin=-48, xmax=48, xtick={-45,-30,-15,0,15,30,45},
  legend pos=north west, legend style={font=\scriptsize},
  grid=major, grid style={gray!20}, tick label style={font=\scriptsize},
  label style={font=\scriptsize}]
  \addplot[thick, blue!70!black, mark=o, mark size=2pt]
    table[x=delta, y=sw] {figs/sw_hw_scan0.dat};
  \addlegendentry{SW (ponto flutuante)}
  \addplot[thick, red, only marks, mark=x, mark size=3pt]
    table[x=delta, y=hw] {figs/sw_hw_scan0.dat};
  \addlegendentry{HW (RTL simulado)}
\end{axis}
\end{tikzpicture}
\caption{Folga por candidato: software \emph{versus} sa\'ida real do RTL, para a
varredura de teste. As curvas se sobrep\~oem em toda a faixa --- o desvio \'e
invis\'ivel nesta escala, o que \'e precisamente o resultado desejado.}
\label{fig:swhw}
\end{figure}

\begin{figure}[!t]
\centering
\begin{tikzpicture}
\begin{axis}[width=\columnwidth, height=4.2cm,
  xlabel={candidato de desvio $\delta_k$ (graus)},
  ylabel={erro HW$-$SW (mm)}, xmin=-48, xmax=48,
  xtick={-45,-30,-15,0,15,30,45},
  legend pos=south west, legend style={font=\tiny},
  grid=major, grid style={gray!20}, tick label style={font=\scriptsize},
  label style={font=\scriptsize}]
  \draw[gray, dashed] (axis cs:-48,0.98) -- (axis cs:48,0.98);
  \draw[gray, dashed] (axis cs:-48,-0.98) -- (axis cs:48,-0.98);
  \node[gray, font=\tiny, anchor=west] at (axis cs:-46,1.7) {$\pm 1$ LSB};
  \addplot[only marks, mark=*, mark size=1.6pt, blue!70!black]
    table[x=delta, y=err_mm] {figs/err_scan0.dat};
  \addlegendentry{varredura 0}
  \addplot[only marks, mark=square*, mark size=1.6pt, red!70!black]
    table[x=delta, y=err_mm] {figs/err_scan1.dat};
  \addlegendentry{varredura 1}
  \addplot[only marks, mark=triangle*, mark size=2pt, orange!90!black]
    table[x=delta, y=err_mm] {figs/err_scan2.dat};
  \addlegendentry{varredura 2}
\end{axis}
\end{tikzpicture}
\caption{Res\'iduo HW$-$SW, com sinal, nas tr\^es varreduras (\num{57}
compara\c{c}\~oes). Linhas tracejadas: $\pm1$ LSB de Q6.10
(\SI{0.98}{\milli\metre}). O erro m\'aximo, \SI{5.27}{\milli\metre}, ocorre quando
a quantiza\c{c}\~ao faz um ponto pr\'oximo da borda do corredor cruzar o limiar
$|y|\le W/2$, trocando qual ponto \'e o m\'inimo --- efeito inerente a um
crit\'erio de decis\~ao por limiar, e n\~ao ac\'umulo de erro aritm\'etico.}
\label{fig:erro}
\end{figure}

\subsection{Desempenho}

O \emph{baseline} de software foi medido na pr\'opria Raspberry~Pi~5 que embarca o
sistema; a lat\^encia de hardware \'e o n\'umero de ciclos \textbf{medido na
simula\c{c}\~ao} (\SI{7515}{\nano\second} a \SI{100}{\mega\hertz}), n\~ao uma
estimativa anal\'itica. A Tab.~\ref{tab:speedup} consolida os tempos.

\begin{table}[!t]
\caption{Tempo por varredura e ganho}
\label{tab:speedup}
\centering\footnotesize
\begin{tabular}{@{}lrrr@{}}
\toprule
\textbf{Implementa\c{c}\~ao} & \textbf{Tempo} & \multicolumn{2}{c}{\textbf{Speedup}} \\
 & & vs.\ escalar & vs.\ NumPy \\
\midrule
Software escalar (Pi 5)      & \SI{3.54}{\milli\second}   & --- & --- \\
Software vetorizado (NumPy)  & \SI{0.488}{\milli\second}  & \num{7.3}$\times$ & --- \\
\textbf{HW \SI{100}{\mega\hertz}} (\num{751} ciclos) & \textbf{\SI{7.51}{\micro\second}} & \num{472}$\times$ & \num{65}$\times$ \\
HW \SI{150}{\mega\hertz}     & \SI{5.01}{\micro\second}   & \num{707}$\times$ & \num{98}$\times$ \\
\bottomrule
\end{tabular}
\end{table}

\begin{figure}[!t]
\centering
\begin{tikzpicture}
\begin{semilogxaxis}[width=\columnwidth, height=4cm,
  xlabel={tempo por varredura ($\mu$s, escala log)},
  xmin=3, xmax=9000, xmajorgrids, grid style={gray!25},
  symbolic y coords={HW150, HW100, NumPy, escalar},
  ytick=data, y dir=reverse,
  yticklabels={HW \SI{150}{\mega\hertz}, HW \SI{100}{\mega\hertz},
               SW NumPy, SW escalar},
  tick label style={font=\scriptsize}, label style={font=\scriptsize},
  nodes near coords, nodes near coords style={font=\tiny, anchor=west},
  point meta=explicit symbolic,
  xbar, bar width=6pt, enlarge y limits=0.25]
  \addplot[fill=red!25, draw=red!60!black] coordinates {
    (3542, escalar) [\SI{3.54}{\milli\second}]
    (488,  NumPy)   [\SI{0.488}{\milli\second}]};
  \addplot[fill=blue!25, draw=blue!60!black] coordinates {
    (7.51, HW100)   [\SI{7.51}{\micro\second}]
    (5.01, HW150)   [\SI{5.01}{\micro\second}]};
\end{semilogxaxis}
\end{tikzpicture}
\caption{Tempo por varredura (escala logar\'itmica --- s\~ao quase tr\^es ordens de
grandeza). \emph{Baseline} de software medido na pr\'opria Raspberry~Pi~5; o
tempo de hardware \'e o n\'umero de ciclos \textbf{medido na simula\c{c}\~ao}.}
\label{fig:tempos}
\end{figure}

Vale interpretar o ganho, e n\~ao apenas exibi-lo (Fig.~\ref{fig:tempos}). O LiDAR
entrega uma varredura a cada \SI{100}{\milli\second}: mesmo o software vetorizado
consome apenas \SI{0.5}{\percent} desse or\c{c}amento, de modo que \textbf{o
acelerador n\~ao ``salva'' o sistema} --- ele n\~ao estava perdendo prazos. O valor real \'e outro,
e triplo: (i)~libera \SI{8.3}{\percent} da CPU num processador que tamb\'em roda
SLAM e EKF, e que j\'a apresentou congelamento por exaust\~ao de mem\'oria;
(ii)~torna a lat\^encia \textbf{determin\'istica}, ao contr\'ario de um n\'o
Python sujeito ao escalonador de um Linux n\~ao-preemptivo --- e determinismo, num
la\c{c}o de seguran\c{c}a, vale mais que m\'edia baixa; (iii)~abre or\c{c}amento
para aumentar $K$ ou a densidade angular sem custo de tempo, j\'a que o gargalo
passa a ser a taxa do sensor.

\subsection{S\'intese: recursos, timing e pot\^encia}

Alvo: \texttt{xc7z020clg484-1} (ZedBoard), per\'iodo de \SI{10}{\nano\second}.
Os valores a seguir vêm dos relat\'orios do Vivado
(\texttt{vivado/reports/}).

\begin{table}[!t]
\caption{S\'intese e implementa\c{c}\~ao (Vivado 2022.2, \texttt{xc7z020clg484-1})}
\label{tab:util}
\centering\footnotesize
\begin{tabular}{@{}lrrrr@{}}
\toprule
\textbf{Recurso} & \textbf{Previsto} & \textbf{Obtido} & \textbf{Dispon\'ivel} & \textbf{\%} \\
\midrule
LUT        & ---          & \num{2445}   & \num{53200}  & \SI{4.6}{\percent} \\
Flip-flop  & ---          & \num{1808}   & \num{106400} & \SI{1.7}{\percent} \\
DSP48E1    & \num{78}     & \textbf{\num{54}} & \num{220} & \SI{24.5}{\percent} \\
BRAM       & \num{0}      & \textbf{\num{0}}  & \num{140} & \SI{0}{\percent} \\
SRL        & ---          & \num{17}     & ---          & --- \\
\midrule
\multicolumn{5}{@{}l@{}}{WNS (\emph{setup}): \textbf{\SI{+0.872}{\nano\second}} $\Rightarrow$ o \emph{timing} fecha a \SI{100}{\mega\hertz}} \\
\multicolumn{5}{@{}l@{}}{WHS (\emph{hold}): \SI{+0.059}{\nano\second}} \\
\multicolumn{5}{@{}l@{}}{$F_{\max}$: \textbf{\SI{109.6}{\mega\hertz}}} \\
\multicolumn{5}{@{}l@{}}{Pot\^encia total: \textbf{\SI{0.246}{\watt}}
  (din\^amica \SI{0.140}{\watt} $+$ est\'atica \SI{0.106}{\watt})} \\
\bottomrule
\end{tabular}
\end{table}

A s\'intese conclui sem erros nem \emph{critical warnings}, e o roteamento fecha
com \textbf{WNS de \SI{+0.872}{\nano\second}} --- folga confort\'avel para o alvo
de \SI{100}{\mega\hertz} ($F_{\max} = \SI{109.6}{\mega\hertz}$). O consumo de
l\'ogica \'e modesto: \SI{4.6}{\percent} dos LUT e \SI{1.7}{\percent} dos
\emph{flip-flops}.

\paragraph{Pot\^encia}
O acelerador consome \SI{0.246}{\watt} no total, dos quais apenas
\SI{0.140}{\watt} s\~ao din\^amicos --- o restante \'e a corrente de fuga do
dispositivo, que existiria com a FPGA ligada mesmo sem nenhum projeto carregado.
A parcela din\^amica \'e baixa porque o \emph{datapath}, embora largo (\num{19}
pistas em paralelo), \'e raso: cada est\'agio faz uma multiplica\c{c}\~ao e uma
compara\c{c}\~ao, sem mem\'oria e sem realimenta\c{c}\~ao. Convém n\~ao
superinterpretar o n\'umero: uma compara\c{c}\~ao honesta com a Raspberry~Pi~5
exigiria medir a pot\^encia \emph{incremental} da CPU ao executar o supervisor, o
que n\~ao foi feito --- e o Zynq tamb\'em consumiria o seu pr\'oprio ARM. O que se
pode afirmar \'e que o custo energ\'etico da l\'ogica \'e desprez\'ivel frente ao
dos motores da cadeira, da ordem de dezenas de \si{\watt}.

\paragraph{Onde a previs\~ao anal\'itica errou --- para mais}
Havia-se previsto \num{78} DSP48 ($19\times4$ das pistas mais \num{2} da
proje\c{c}\~ao). A ferramenta usou \textbf{\num{54}}, \SI{31}{\percent} a menos.
A raz\~ao \'e instrutiva e s\'o aparece ao ler o \emph{DSP Final Report}: o
DSP48E1 n\~ao \'e um multiplicador isolado, mas um bloco
\emph{multiplicador-acumulador} com uma porta $C$ e uma cascata dedicada
(\texttt{PCIN}) para o pr\'oximo DSP da coluna. A s\'intese reconheceu que
\begin{equation*}
x^{(k)} = X\cos\delta_k + Y\sin\delta_k
\end{equation*}
\noindent \'e exatamente a forma que o bloco implementa nativamente, e fundiu o
segundo produto e a soma em \emph{um} DSP --- ora pela porta $C$
(\texttt{C+A*B}), ora encadeando pelo \texttt{PCIN} (\texttt{PCIN+A*B}). Onde o
coeficiente \'e trivial, economizou ainda mais: a pista central ($\delta_9=0$,
com $\sin\delta_9 = 0$) dispensa metade dos produtos.

A li\c{c}\~ao \'e que estimar recursos contando operadores do RTL \emph{superestima}
o custo: o mapeamento explora estrutura do dispositivo que n\~ao aparece no
c\'odigo. A folga resultante \'e substancial --- sobram \SI{75}{\percent} dos DSP,
espa\c{c}o para praticamente \emph{triplicar} o n\'umero de candidatos $K$ sem
sair do \texttt{xc7z020}.

\paragraph{Confirma\c{c}\~oes} O CORDIC de fato n\~ao consome
\textbf{nenhum} DSP --- todos os \num{54} est\~ao no \texttt{corridor\_core}
(pistas e proje\c{c}\~ao), como se pretendia ao escolher somadores e
deslocamentos cabeados. E o uso de BRAM \'e \textbf{zero}, confirmando que a
arquitetura \emph{streaming} n\~ao precisa armazenar a varredura: cada feixe \'e
consumido no ciclo em que chega.

% =============================================================================
\section{Conclus\~ao}

O acelerador atinge o objetivo: \num{472}$\times$ mais r\'apido que o software
escalar e \num{65}$\times$ que o vetorizado, com equival\^encia \textbf{bit a
bit} verificada contra um modelo de refer\^encia sobre dados reais da cadeira.

A an\'alise cr\'itica, por\'em, precisa ir al\'em do n\'umero. \textbf{O ganho de
tempo n\~ao era necess\'ario}: o software j\'a cabia folgadamente no or\c{c}amento
de \SI{100}{\milli\second} do sensor. O benef\'icio leg\'itimo \'e o alívio de CPU
numa Pi que roda SLAM e EKF, e sobretudo o \textbf{determinismo} da lat\^encia num
la\c{c}o de seguran\c{c}a. Reportamos isso explicitamente porque a conclus\~ao
oposta --- ``aceleramos, logo o sistema melhorou'' --- seria uma leitura
desonesta dos dados.

O resultado mais transfer\'ivel n\~ao \'e o \emph{speedup}, e sim o
\textbf{defeito de faixa em Q3.13}: um \^angulo de \SI{-5.85}{\radian} saturando
num formato de $\pm\SI{4}{\radian}$ produziu \SI{402}{\milli\metre} de erro
\emph{silencioso} --- o RTL n\~ao acusa, a simula\c{c}\~ao ``roda'', e s\'o a
compara\c{c}\~ao com uma refer\^encia f\'isica revela. Foi o modelo de ouro, e
n\~ao a inspe\c{c}\~ao de ondas, que o encontrou.

\subsection{Trabalhos futuros}
\begin{enumerate}
\item \textbf{Gravar em placa}: os resultados aqui s\~ao de simula\c{c}\~ao
p\'os-s\'intese; falta o ensaio no ZedBoard, com o driver do ARM
(\texttt{sw/lidar\_accel.c}) sobre o AXI-DMA.
\item \textbf{Explorar o espa\c{c}o de projeto}: com \SI{35}{\percent} dos DSP
usados, cabem $K$ maior (mais candidatos) ou pistas com maior precis\~ao;
resta caracterizar o \emph{trade-off} recursos $\times$ pot\^encia $\times$
qualidade da decis\~ao.
\item \textbf{Reduzir a lat\^encia fim-a-fim}: o acelerador resolve
\SI{7.5}{\micro\second} de um la\c{c}o cujo verdadeiro atraso \'e a
\SI{100}{\milli\second} do sensor. O ganho maior estaria em processar a varredura
\emph{incrementalmente}, \`a medida que os feixes chegam --- algo que a
arquitetura \emph{streaming} j\'a permite, e o software n\~ao.
\end{enumerate}

% =============================================================================
\begin{thebibliography}{9}

\bibitem{fehr2000}
L.~Fehr, W.~E. Langbein e S.~B. Skaar, ``Adequacy of power wheelchair control
interfaces for persons with severe disabilities: A clinical survey,''
\emph{J. Rehabil. Res. Dev.}, vol.~37, no.~3, pp.~353--360, 2000.

\bibitem{borenstein1991}
J.~Borenstein e Y.~Koren, ``The Vector Field Histogram --- Fast Obstacle
Avoidance for Mobile Robots,'' \emph{IEEE Trans. Robot. Autom.}, vol.~7, no.~3,
pp.~278--288, 1991.

\bibitem{volder1959}
J.~E. Volder, ``The CORDIC Trigonometric Computing Technique,'' \emph{IRE Trans.
Electron. Comput.}, vol.~EC-8, no.~3, pp.~330--334, 1959.

\bibitem{andraka1998}
R.~Andraka, ``A survey of CORDIC algorithms for FPGA based computers,'' em
\emph{Proc. ACM/SIGDA Int. Symp. FPGAs}, 1998, pp.~191--200.

\bibitem{zynqbook}
L.~H. Crockett, R.~A. Elliot, M.~A. Enderwitz e R.~W. Stewart, \emph{The Zynq
Book: Embedded Processing with the ARM Cortex-A9 on the Xilinx Zynq-7000}.
Glasgow: Strathclyde Academic Media, 2014.

\bibitem{ros2docs}
Open Robotics, ``ROS\,2 Jazzy Jalisco Documentation,'' 2024. [Online].
Dispon\'ivel: \url{https://docs.ros.org/en/jazzy/}

\end{thebibliography}

\end{document}
